\chapter{Kva eit fag er}

\begin{motto}
Eit fundament er ikkje ein formel alle deler.\\ Det er ein måte å vere usamd på som fører ein stad.
\end{motto}

\textsc{Grunnlaust} har eit underforstått bilete av kva eit fag med eit fundament ser ut som, og heile saka mot design heng på det biletet. I biletet har det modne faget éin grunnstruktur som alle deler — ein formel, ein teori, ein metode — og denne strukturen gjer at to fagfolk uavhengig kjem til same dom av same grunn. Fysikken er førebiletet. Dette kapitlet bestrider biletet. Det viser at svært få modne fag faktisk ser slik ut, at fundamentet til eit fag ligg ein annan stad enn i ein delt formel, og at design har eit fundament i nett den forstanden som faktisk gjeld for profesjonsfag.

\section{Dei pluralistiske faga}

Tenk på kva slags fag verda faktisk reknar som modne og velgrunna. Jussen er eit. Han har inga formel som avgjer dommar; to dyktige juristar kan lese same lova og same faktum og komme til motsette konklusjonar, kvar med ei grunngjeving fagfellane må respektere. Likevel tvilar ingen på at jussen er ein disiplin med eit fundament. Klinisk medisin er eit anna. I møtet med den einskilde pasienten er medisinen full av skjøn som ikkje let seg redusere til protokoll; røynde legar er usamde om diagnose og behandling dagleg, og lærebøkene innrømmer det. Psykiatrien, pedagogikken, historiefaget, sosialantropologien, til og med store delar av økonomien: alle er fag der skjønet sit i kjernen, der usemje mellom kompetente utøvarar er normaltilstanden, og der det ikkje finst ein delt formel som tvingar fram éin dom.

Poenget er ikkje at desse faga er dårlege fag. Poenget er at \emph{fundamentet deira ligg ikkje i ein delt formel}. Det ligg ein annan stad: i eit delt sett omgrep, ein delt litteratur, ein delt utdanningsveg, ein delt praksis for å føre og prøve argument, og ein delt dom over kva som tel som eit godt og eit dårleg arbeid sjølv når utfallet er omstridd. Eit fag har eit fundament når usemja i det er \emph{disiplinert} — når ho følgjer reglar, byggjer på felles referansar, og fører ein stad — ikkje når ho er avskaffa. Den som krev at design skal ha eit fundament av fysikken sitt slag, har ikkje sett ein høg standard for design. Han har sett ein standard nesten ingen fag når, og så peika ut design som om mangelen var særeigen.

\section{Schöns korreksjon av det vitskaplege biletet}

Det var Donald Schön som mest grundig avkledde biletet av at profesjonell kunnskap er bruk av vitskap \autocite{schon1983}. Han kalla biletet \emph{teknisk rasjonalitet}: ideen om at den dugande profesjonelle fyrst lærer ein vitskapleg teori og så bruker han på praktiske problem. Schön viste, gjennom nære studiar av korleis profesjonelle faktisk arbeider, at dette ikkje skildrar nokon profesjon — verken ingeniøren, legen, advokaten eller arkitekten. Den verkelege kompetansen ligg i ein refleksjon-i-handling som ikkje er mindre rigorøs enn teori, men annleis: ein dialog med situasjonen der utøvaren rammar problemet, prøver eit grep, les responsen, og rammar på nytt. Dette er ikkje fråvêret av eit fundament. Det er eit fundament av eit anna slag enn det den tekniske rasjonaliteten føresette.

Schöns poeng rammar \textsc{Grunnlaust} direkte. Når kritikaren seier at design manglar eit fundament fordi det ikkje kan grunngje dommane sine slik fysikken kan, måler han design mot teknisk rasjonalitet — det same biletet Schön viste at ikkje eingong dei harde profesjonane lever opp til. Ingeniøren grunngjev ikkje den faktiske ingeniørgjerninga si med rein mekanikk; han bruker skjøn, røynsle, tommelfingerreglar, og ein refleksjon-i-handling som er like «privat» som designaren sin, heilt til han skriv han ned. At designaren sin kunnskap er taus, gjer han ikkje verdlaus eller uoverførbar; han overførast slik all profesjonskunnskap overførast, gjennom opplæring i ein praksis.

\section{Kva design faktisk deler}

Sett at vi sluttar å krevje ein formel og i staden spør om design har det dei pluralistiske faga har. Då ser biletet heilt annleis ut. Design har eit delt sett omgrep — fra affordans og brukar til iterasjon, prototype og innramming — som tyder noko nær det same for kompetente utøvarar verda over. Det har ein delt litteratur, med kanoniske verk og fagfellevurderte tidsskrift. Det har ein delt utdanningsveg med atelier, kritikk og portefølje. Det har ein delt praksis for å føre og prøve arbeid, og ein dom — om enn omstridd i einskildtilfelle — over kva som skil eit gjennomarbeidd prosjekt frå eit slurvete. Med andre ord har design nett den typen fundament dei andre profesjonsfaga har: ikkje ein formel, men ein disiplinert praksis med delte referansar.

At dette fundamentet ikkje liknar fysikken sitt, gjer det ikkje til eit ikkje-fundament. Det gjer det til eit fundament av det slaget profesjonsfag har. Simon såg dette då han plasserte design saman med medisin, ingeniørfag og jus som «vitskapane om det kunstige» \autocite{simon1969sciences} — fag som ikkje handlar om korleis verda \emph{er}, men om korleis ho \emph{kan bli}, og som difor har ein annan logikk enn naturvitskapen. Og Cross bygde dette ut til ein heil epistemologi for design som disiplin \autocite{cross2006}. Spørsmålet er ikkje om design har eit fundament. Spørsmålet er om vi er viljuge til å sjå eit fundament som ikkje har forma vi venta.

\section{Det før-paradigmatiske, igjen}

Eitt motargument står att, og det er Kuhn-argumentet frå \textsc{Grunnlaust}: at design er \emph{før-paradigmatisk}, eit fag før sitt fyrste delte rammeverk, kjenneteikna av skular som ikkje kan avgjere usemjene sine \autocite{kuhn1962}. Men her må vi vere varsame med å bruke Kuhn mot design, av to grunnar. For det fyrste skreiv Kuhn om naturvitskapane; han hevda aldri at samfunnsfag og profesjonsfag måtte gjennom same paradigmestruktur, og mange av dei har det openbert ikkje, utan at vi kallar dei avleggarar. For det andre er det eit ope empirisk spørsmål, ikkje ein avgjord sak, om design har konkurrerande «paradigme» eller om det har eit modnande, felles rammeverk som veks seg sterkare — slik den samtidige designforskinga, og særleg den nordiske, faktisk hevdar at det gjer. Dei følgjande delane legg fram det rammeverket: kunnskapsforma i \textsc{del ii}, metoden og kumulasjonen i \textsc{iii}, og dei norske bidraga i \textsc{vi}. Lesaren får sjølv avgjere om dette er eit fag på veg \emph{mot} eit paradigme — altså eit ungt fag, ikkje ein avleggar — eller om kritikaren har rett. Men spørsmålet kan ikkje avgjerast ved definisjon. Det må avgjerast ved å sjå kva faget faktisk har bygd.

\vignett

Og det fyrste det har bygd, er ei eiga kunnskapsform. Det er emnet for neste del.

\vidarelesing{\fullcite{schon1983}; \fullcite{cross2006}; \fullcite{simon1969sciences}; \fullcite{kuhn1962}.}
